\begin{recipebox}{Sticky Date Pudding}{8}{1 hr}
\begin{ingredients}
  \item[] \textbf{Pudding}
  \item 1\nicefrac{1}{4} cups pitted dates, chopped
  \item 1\nicefrac{1}{4} cups boiling water
  \item 1 tsp baking soda
  \item 60\,g butter, softened
  \item \nicefrac{3}{4} cup brown sugar
  \item 2 eggs
  \item 1 cup self-raising flour
  \item[] \textbf{Butterscotch sauce}
  \item 2 cups brown sugar
  \item 250\,ml cream
  \item 50\,g butter
  \item 1 tsp vanilla essence
\end{ingredients}

\begin{method}
  \item Preheat oven to 180\degree C bake (160\degree C fan). Grease and line a 20cm square cake tin.
  \item Put the dates in a bowl, pour over the boiling water and stir in the baking soda (it will foam). Leave for 5 minutes.
  \item Cream the butter and brown sugar, then beat in the eggs one at a time.
  \item Fold in the flour, then the dates and all their liquid. The batter will be quite runny.
  \item Pour into the tin and bake for 35--40 minutes, until a skewer comes out clean.
  \item Meanwhile, make the sauce: stir the brown sugar, cream, butter and vanilla in a saucepan over low heat until the sugar dissolves and the butter melts, then simmer for 2--3 minutes.
  \item Serve the pudding warm in squares with the sauce poured over.
\end{method}

\begin{notes}
Quantities follow Chelsea Sugar's sticky date pudding with butterscotch sauce. The sauce makes plenty --- leftovers keep in the fridge for a week.
The pudding freezes well in portions. Reheat a square in the microwave and pour over warm sauce.
\end{notes}
\end{recipebox}
